\chapter{Solution to exercises in chapter 3}
\begin{enumerate}
\item\label{c3p1} Matrix element of Fock operator.

The Fock operator is
\[
\begin{aligned}
f(1) &= h(1) + \sum_b\left(\int d\vec{x}_2|\chi_b(\vec{x}_2)|^2 r_{12}^{-1} -
\int d\vec{x}_2\chi_b^\ast(\vec{x}_2)r_{12}^{-1}P_{12}\chi_b(\vec{x}_2)\right) \\
 &= h(1) + \sum_b\left(
\int d\vec{x}_2\chi_b^\ast(\vec{x}_2)r_{12}^{-1}\chi_b(\vec{x}_2) -
\int d\vec{x}_2\chi_b^\ast(\vec{x}_2)r_{12}^{-1}P_{12}\chi_b(\vec{x}_2)\right) 
\end{aligned}
\]
so that
\[
\begin{aligned}
f\ket{j} &= h(\vec{x}_1)\chi_j(\vec{x}_1) + \\
 & \sum_{b \ne j}\left(\int d\vec{x}_2
 \chi_b^\ast(\vec{x}_2)r_{12}^{-1}\chi_b(\vec{x}_2)\chi_j(\vec{x}_1) -
\int d\vec{x}_2\chi_b^\ast(\vec{x}_2)r_{12}^{-1}P_{12}\chi_b(\vec{x}_2)
\chi_i(\vec{x}_1) \right) 
\end{aligned}
\]
and
\[
\begin{aligned}
\braket{i|f|j} &= \braket{i|h|j} + \\
 & \sum_{b \ne j}\left(\iint d\vec{x}_1 d\vec{x}_2 \chi_i^\ast(\vec{x}_1)
 \chi_b^\ast(\vec{x}_2)r_{12}^{-1}\chi_b(\vec{x}_2)\chi_j(\vec{x}_1)\right. - \\
 & \left.
 \iint d\vec{x}_1 d\vec{x}_2
 \chi_i^\ast(\vec{x}_1)\chi_b^\ast(\vec{x}_2)r_{12}^{-1}P_{12}\chi_b(\vec{x}_2)
 \chi_j(\vec{x}_1) \right) \\
 &= \braket{i|h|j} + \\
 & \sum_{b \ne j}\left(\iint d\vec{x}_1 d\vec{x}_2 \chi_i^\ast(\vec{x}_1)
\chi_j(\vec{x}_1) r_{12}^{-1}\chi_b(\vec{x}_2)\chi_b^\ast(\vec{x}_2)\right. - \\
 & \left.
 \iint d\vec{x}_1 d\vec{x}_2
 \chi_i^\ast(\vec{x}_1)\chi_b^\ast(\vec{x}_2)r_{12}^{-1}\chi_b(\vec{x}_1)
 \chi_j(\vec{x}_2) \right) \\
 &= \braket{i|h|j} + \\
 & \sum_{b \ne j}\left(\iint d\vec{x}_1 d\vec{x}_2 \chi_i^\ast(\vec{x}_1)
\chi_j(\vec{x}_1) r_{12}^{-1}\chi_b(\vec{x}_2)\chi_b^\ast(\vec{x}_2)\right. - \\
 & \left.
 \iint d\vec{x}_1 d\vec{x}_2
 \chi_i^\ast(\vec{x}_1)\chi_b(\vec{x}_1)r_{12}^{-1}\chi_b^\ast(\vec{x}_2)
 \chi_j(\vec{x}_2) \right) \\
 &= \braket{i|h|j} + \sum_b\left([ij|bb] - [ib|bj]\right)
\end{aligned}
\]

\item\label{c3p2} Lagrange multipliers are real.
Equation \eqref{c3e17} is
\[
\mathscr{L}[\{\chi_a\}] = \braket{\Psi_0|H|\Psi_0} - \sum_{ab}\varepsilon_{ab}(
[a|b] - \delta_{ab})
\]
Since $H$ is hermitian, $\braket{\Psi_0|H|\Psi_0}$ is real. Because of 
orthonormality, $[a|b] = \braket{\chi_a|\chi_b}$ is real. $\delta_{ab}$ is also
a real number. The functional $\mathscr{L}$ seeks to minimise energy, a real
quantity. Therefore, it has to be real. As a result, $\varepsilon_{ab}$ too
must be real.

\item\label{c3p3} Calculation of $\delta E_0$.
Done in the notes to arrive at equation \eqref{c3e25}.

\item Hermiticity of the Fock operator.
We showed in problem \ref{c3p1} that
\[
\braket{i|f|j} = \braket{i|h|j} + \sum_b\left([ij|bb] - [ib|bj]\right)
\]
so that
\[
\begin{aligned}
\braket{j|f|i} &= \braket{j|h|i} + \sum_b\left([ji|bb] - [jb|bi]\right) \\
 &= \braket{i|h|j}^\ast + \sum_b\left([ij|bb]^\ast - [jb|bi]\right)
\end{aligned}
\]
The first term changed its form because of hermiticity of the one-electron
hamiltonian in the field of nuclei. The term $[jb|bi]$ needs a bit of care. 
From equation \eqref{c2e90}
\[
[ij|kl] = \iint d\vec{x}_1d\vec{x}_2\chi_i^\ast(\vec{x}_1)\chi_j(\vec{x}_1)
r_{12}^{-1}\chi_k^\ast(\vec{x}_2)\chi_l(\vec{x}_2).
\]
so that
\[
\begin{aligned}
[jb|bi] &= \iint d\vec{x}_1d\vec{x}_2\chi_j^\ast(\vec{x}_1)\chi_b(\vec{x}_1)
r_{12}^{-1}\chi_b^\ast(\vec{x}_2)\chi_i(\vec{x}_2) \\
&= \left[\iint d\vec{x}_1d\vec{x}_2\chi_j(\vec{x}_1)\chi_b^\ast(\vec{x}_1)
r_{12}^{-1}\chi_b(\vec{x}_2)\chi_i^\ast(\vec{x}_2) \right]^\ast \\
&= \left[\iint d\vec{x}_1d\vec{x}_2\chi_i^\ast(\vec{x}_2)\chi_b(\vec{x}_2)
r_{12}^{-1}\chi_b^\ast(\vec{x}_1)\chi_j(\vec{x}_1) \right]^\ast \\
&= \left[\iint d\vec{x}_1d\vec{x}_2\chi_i^\ast(\vec{x}_1)\chi_b(\vec{x}_1)
r_{12}^{-1}\chi_b^\ast(\vec{x}_2)\chi_j(\vec{x}_2) \right]^\ast \\
&= [ib|bj]^\ast
\end{aligned}
\]
so that
\[
\begin{aligned}
\braket{j|f|i}&=\braket{i|h|j}^\ast+\sum_b\left([ij|bb] - [jb|bi]\right)^\ast\\
&= \braket{i|f|j}^\ast
\end{aligned}
\]
making $f$ a hermitial operator.

\item\label{c3p5} Energy needed to eject two electrons.
\[
\begin{aligned}
^NE_0 &= \sum_{a=1}^N\braket{a|h|a} + \frac{1}{2}\sum_{a,b=1}^N\braket{ab||ba} \\
^{N-2}E_{cd} &= \sum_{a=1, a\ne c,d}^N\braket{a|h|a} +
\frac{1}{2}\sum_{a\ne c,d}\sum_{b \ne c,d}\braket{ab||ba}
\end{aligned}
\]
so that
\[
\begin{aligned}
^{N-2}E_{c,d} - ^NE_0 &= -\braket{c|h|c} - \braket{d|h|d} - \\
 &\quad \frac{1}{2}\sum_{b=1}^N\braket{cb||bc} - 
   \frac{1}{2}\sum_{b=1}^N\braket{db||bd} - \\
 &\quad \frac{1}{2}\sum_{a=1}^N\braket{ac||ca} - 
  \frac{1}{2}\sum_{a=1}^N\braket{ad||da}.
\end{aligned}
\]
Using \eqref{c2e86} and \eqref{c2e88},
\[
\begin{aligned}
^{N-2}E_{c,d} - ^NE_0 &= -\braket{c|h|c} - \braket{d|h|d} -
\sum_{b=1}^N\braket{cb||bc} - \sum_{b=1}^N\braket{db||bd} \\
 &= -\varepsilon_c - \varepsilon_d
\end{aligned}
\]

\item\label{c3p6} Electron affinity. Calculated as \eqref{c3e48}.

\item\label{c3p7} Commutator of Hartree-Fock hamiltonian and the permutation
operator.

Consider the general permutation operator $P_{ij}$ which swaps electrons $i$
and $j$. It is an $N$-electron operator that acts on $N$-electron
wavefunctions. On the other hand, $f(1)$ is a one electron operator. Therefore,
it does not make sense to check if they commute. We are asked to investigate
$[P_{ij}, H]$, where $H$ is the Hartree-Fock hamiltonian which is the sum of
the Fock operators for all electrons. We want to show that
\[
\left[P_{ij}, \sum_k f(k)\right] = 0.
\]
Without loss of generality, we will choose $i = 1, j = 2$. Then,
\[
P_{12}\Psi(\vec{x}_1, \vec{x}_2, \ldots, \vec{x}_N) = \Psi(\vec{x}_2,
\vec{x}_1, \ldots, \vec{x}_N).
\]
For $k \ge 3$,
\[
P_{12}f(k)\Psi(\vec{x}_1, \vec{x}_2, \ldots, \vec{x}_N) = f(k)\Psi(\vec{x}_2,
\vec{x}_1, \ldots, \vec{x}_N) = f(k)P_{12}\Psi(\vec{x}_1, \vec{x}_2, \ldots,
\vec{x}_N),
\]
that is $[P_{12}, f(k)] = 0$ for $k \ge 3$.

Now consider
\[
\begin{aligned}
f(1)\Psi(\vec{x}_1, \ldots, \vec{x}_N) &= h(1)\Psi(\vec{x}_1, \ldots,
\vec{x}_N) + \\
 &= \sum_b\int \frac{d\vec{x}^\op}{|\vec{r}_1 - \vec{r}^\op|}\chi_b^\ast(\vec{x}^\op)
 (\chi_b(\vec{x}^\op) - \chi_b(\vec{x}_1))\Psi(\vec{x}_1, \ldots, \vec{x}_N).
\end{aligned}
\]
so that
\[
\begin{aligned}
P_{12}f(1)\Psi(\vec{x}_1, \ldots, \vec{x}_N) &= h(2)\Psi(\vec{x}_2, \vec{x}_1,
\ldots, \vec{x}_N) + \\
 & \sum_b\int\frac{d\vec{x}^\op}{|\vec{r}_2 -
 \vec{r}^\op|}\chi_b^\ast(\vec{x}^\op)(\chi_b(\vec{x}^\op) - \chi_b(\vec{x}_2))
 \Psi(\vec{x}_2, \vec{x}_1, \ldots, \vec{x}_N) \\
 &= h(2)P_{12}\Psi(\vec{x}_1, \vec{x}_2, \ldots, \vec{x}_N) + \\
 & \sum_b\int\frac{d\vec{x}^\op}{|\vec{r}_2 -
 \vec{r}^\op|}\chi_b^\ast(\vec{x}^\op)(\chi_b(\vec{x}^\op) - \chi_b(\vec{x}_2))
 P_{12}\Psi(\vec{x}_1, \vec{x}_2, \ldots, \vec{x}_N) \\
 &= f(2)P_{12}\Psi(\vec{x}_1, \vec{x}_2, \ldots, \vec{x}_N) 
\end{aligned}
\]
Thus, $P_{12}f(1) = f(2)P_{12}$. From this, we also get $P_{12}^2 f(1) =
P_{12}f(2)P_{12} \Rightrrow f(1) = P_{12}f(2)P_{12} \Rightarrow f(1)P_{12} =
P_{12}f(2)P_{12}^2 \Rightrrow f(1)P_{12} = P_{12}f(2)$.

Therefore,
\[
\begin{aligned}
P_{12}(f(1) + f(2) + f(3) + \cdots + f(N)) &= f(2)P_{12} + f(1)P_{12} +
f(3)P_{12} + \cdots f(N)P_{12} \\
&= (f(1) + \cdots f(N))P_{12}.
\end{aligned}
\]
Thus, $P_{12}$ commutes with the Hartree-Fock hamiltonian.
\end{enumerate}
